\chapter{Funkcje elementarne}
\label{ma1:ch:elementary-functions}

Funkcje wykładniczą i trygonometryczne skonstruujemy za pomocą szeregów liczbowych. Ich prawa algebraiczne będą konsekwencjami iloczynu Cauchy'ego, a nie założeniami definicyjnymi. Nie korzystamy z pochodnych, całek ani geometrycznej definicji liczby $\pi$.

\section{Funkcja wykładnicza i logarytm}
\begin{definicja}
\label{ma1:def:exponential-series}
Dla $z\in\C$ określamy $\exp z=\sum_{n=0}^\infty z^n/n!$, przyjmując $0!=1$ i $z^0=1$.
\end{definicja}
Szereg jest bezwzględnie zbieżny: dla $|z|\le R$, $R>0$, oraz $n+1\ge2R$ stosunek kolejnych dodatnich majorant $R^n/n!$ jest nie większy niż $1/2$. Ogon majorant jest więc ograniczony przez ogon odpowiedniego szeregu geometrycznego.

\begin{twierdzenie}
\label{ma1:thm:exp-addition}
Funkcja $\exp:\C\to\C$ jest ciągła, $\exp 0=1$, $\exp(z+w)=\exp z\exp w$ i $\exp z\ne0$ dla wszystkich $z,w\in\C$.
\end{twierdzenie}
\begin{proof}
Dla $|z|,|w|\le R$ i $n\ge1$ rozkład różnicy potęg daje
$|z^n-w^n|\le nR^{n-1}|z-w|$. Stąd, po przejściu do granicy sum częściowych,
\[
 |\exp z-\exp w|\le |z-w|\sum_{n=1}^\infty\frac{R^{n-1}}{(n-1)!}
 =\exp(R)|z-w|.
\]
Każdy punkt ma otoczenie zawarte w pewnym takim kole, co daje ciągłość. Z iloczynu Cauchy'ego szeregów bezwzględnie zbieżnych oraz wzoru dwumianowego otrzymujemy
\[
 \exp z\exp w=\sum_{n=0}^\infty\sum_{k=0}^n
 \frac{z^kw^{n-k}}{k!(n-k)!}
 =\sum_{n=0}^\infty\frac{(z+w)^n}{n!}=\exp(z+w).
\]
Definicja daje $\exp 0=1$, zatem $\exp z\exp(-z)=1$, co dowodzi niezerowości.
\end{proof}
\begin{twierdzenie}
\label{ma1:thm:real-exp-log}
Ograniczenie $\exp$ do $\R$ jest ściśle rosnącą ciągłą bijekcją $\R\to(0,+\infty)$. Jej odwrotność, oznaczana przez $\log$, jest ciągła, ściśle rosnąca i spełnia
$\log(uv)=\log u+\log v$ dla $u,v>0$.
\end{twierdzenie}
\begin{proof}
Dla rzeczywistego $x$ mamy $\exp x=(\exp(x/2))^2>0$. Dla $t>0$ wszystkie składniki szeregu są nieujemne, więc $\exp t\ge1+t>1$. Jeśli $x<y$, to $\exp y=\exp x\exp(y-x)>\exp x$. Ponadto $\exp t\ge1+t$ dla $t\ge0$, zatem $\exp t\to+\infty$ przy $t\to+\infty$, a z tożsamości $\exp(-t)=1/\exp t$ otrzymujemy granicę $0$ w $-\infty$. Twierdzenie o wartości pośredniej daje każdą dodatnią wartość. Ścisła monotoniczność zapewnia jednoznaczność. Ciągłość odwrotności wynika z twierdzenia \ref{ma1:thm:strict-monotone-inverse}. Dla $u,v>0$ podstawiamy $x=\log u$, $y=\log v$ w prawie dodawania argumentów i stosujemy $\log$.
\end{proof}
Niech $e=\exp 1$. Mamy $\log1=0$, $\log e=1$, $\log(1/u)=-\log u$. Logarytm dąży do $+\infty$ przy $u\to+\infty$ i do $-\infty$ przy $u\to0^+$: dla dowolnego $M>0$ wystarczy odpowiednio $u>\exp M$ lub $0<u<\exp(-M)$.

\section{Potęgi rzeczywiste}
\begin{definicja}
\label{ma1:def:real-powers}
Dla $a>0$, $x\in\R$ przyjmujemy $a^x=\exp(x\log a)$. Dla $a>0$, $a\ne1$ definiujemy $\log_a u=\log u/\log a$, $u>0$.
\end{definicja}
\begin{twierdzenie}
\label{ma1:thm:power-laws}
Dla $a,b>0$ i $x,y\in\R$ zachodzą $a^{x+y}=a^xa^y$, $(ab)^x=a^xb^x$, $(a^x)^y=a^{xy}$ oraz $a^0=1$. Funkcja $x\mapsto a^x$ jest ciągła; dla $a>1$ jest ściśle rosnąca, dla $0<a<1$ ściśle malejąca, a dla $a=1$ stała. Dla $a\ne1$ jej odwrotnością na $(0,+\infty)$ jest $\log_a$.
\end{twierdzenie}
\begin{proof}
Pierwsza równość wynika z prawa dodawania argumentów funkcji $\exp$, druga z prawa logarytmowania iloczynu. Ponieważ $\log(a^x)=x\log a$, otrzymujemy trzecią równość. Pozostałe tezy wynikają z $\exp0=1$, ciągłości złożenia, znaku $\log a$ i bijektywności $\exp$.
\end{proof}
Nowa definicja jest zgodna z wcześniejszą konstrukcją z lematu \ref{ma1:lem:holder-powers}. Dla $m\in\Z$, $n\ge1$ liczba $\exp((m/n)\log a)$ jest dodatnia, a jej $n$-ta potęga wynosi $a^m$ (dla ujemnego $m$ przez odwrotność). Jednoznaczność dodatniego pierwiastka daje zgodność dla wykładników wymiernych. Dla $0<\alpha\le1$ i ciągu wymiernego z poprzedniej konstrukcji ciągłość $\exp$ daje tę samą granicę. Dla $x>0$ można przyjąć $0^x=0$; przy ustalonym $x>0$ funkcja $a\mapsto a^x$ jest ciągła również w zerze, ponieważ $\log a\to-\infty$ i $\exp(x\log a)\to0$. Nie przypisujemy w tej definicji wartości $0^0$ ani $0^x$ dla $x<0$.

\section{Sinus, cosinus i wzory dodawania}
\begin{definicja}
\label{ma1:def:trig-series}
Dla $t\in\R$ określamy
\[
 \cos t=\sum_{n=0}^\infty\frac{(-1)^nt^{2n}}{(2n)!},
 \qquad
 \sin t=\sum_{n=0}^\infty\frac{(-1)^nt^{2n+1}}{(2n+1)!}.
\]
\end{definicja}
Bezwzględna zbieżność wynika z majoranty szeregu wykładniczego. Rozdzielenie wyrazów parzystych i nieparzystych w bezwzględnie zbieżnym szeregu daje wzór Eulera $\exp(it)=\cos t+i\sin t$. Z ciągłości $\exp$ oraz nierówności $|\operatorname{Re}z|,|\operatorname{Im}z|\le|z|$ wynika ciągłość obu funkcji. Bezpośrednio z definicji $\cos$ jest parzysta, $\sin$ nieparzysta, $\cos0=1$, $\sin0=0$.

\begin{twierdzenie}
\label{ma1:thm:trig-addition}
Dla $s,t\in\R$ mamy
\begin{align*}
 \cos(s+t)&=\cos s\cos t-\sin s\sin t,\\
 \sin(s+t)&=\sin s\cos t+\cos s\sin t,
 \qquad \cos^2t+\sin^2t=1.
\end{align*}
\end{twierdzenie}
\begin{proof}
Stosujemy prawo dodawania argumentów $\exp$ do $is,it$ i porównujemy części rzeczywiste i urojone iloczynu. Podstawiając $s=-t$ we wzorze dla cosinusa i korzystając z parzystości oraz nieparzystości, otrzymujemy ostatnią równość.
\end{proof}
W szczególności $|\sin t|,|\cos t|\le1$. Wzory dla różnicy argumentów wynikają z zastąpienia $t$ przez $-t$. Podstawiając $s=t$, otrzymujemy $\sin(2t)=2\sin t\cos t$ i $\cos(2t)=\cos^2t-\sin^2t=1-2\sin^2t$.

\section{Liczba $\pi$, monotoniczność i okresowość}
\begin{twierdzenie}
\label{ma1:thm:pi-construction}
Istnieje najmniejsze dodatnie zero $a$ funkcji $\cos$, przy czym $0<a<2$, $\sin a=1$. Definiujemy $\pi=2a$. Wówczas $\cos\pi=-1$, $\sin\pi=0$, a obie funkcje mają okres $2\pi$.
\end{twierdzenie}
\begin{proof}
Ciągłość i $\cos0=1$ dają $\rho>0$, dla którego $\cos t>0$ przy $0\le t\le\rho$. Szereg naprzemienny dla $\cos2$, od wyrazu o indeksie $n=1$, ma malejące moduły, ponieważ $4/((2n+1)(2n+2))<1$. Z oszacowania reszty otrzymujemy $\cos2\le1-2+2/3=-1/3$. Twierdzenie o wartości pośredniej daje zero w $(0,2)$. Niech $a$ będzie infimum zbioru zer w $[0,2]$. Wybierając zera z przedziałów $[a,a+1/n)$ i stosując ciągłość, dostajemy $\cos a=0$ oraz $0<a<2$. Dla $0\le t<a$ mamy $\cos t>0$, gdyż wartość niedodatnia dawałaby wcześniejsze zero przez twierdzenie o wartości pośredniej.

Dla $0<t\le2$ moduły kolejnych wyrazów szeregu sinus maleją, gdyż $t^2/((2n+2)(2n+3))\le4/6<1$. Stąd $\sin t\ge t-t^3/6>0$. W szczególności $\sin a>0$, więc z tożsamości kwadratów $\sin a=1$. Wzory podwojenia dają $\cos(2a)=-1$ i $\sin(2a)=0$, a ponownie $\cos(4a)=1$ i $\sin(4a)=0$. Wzory dodawania argumentów dowodzą okresowości.
\end{proof}
\begin{twierdzenie}
\label{ma1:thm:trig-monotonicity}
Sinus jest ściśle rosnący na $[-\pi/2,\pi/2]$ i przyjmuje tam wszystkie wartości z $[-1,1]$. Cosinus jest ściśle malejący na $[0,\pi]$ i przyjmuje tam wszystkie wartości z $[-1,1]$. Najmniejszy dodatni okres każdej z tych funkcji wynosi $2\pi$.
\end{twierdzenie}
\begin{proof}
Dla $0<t<\pi$ mamy $\sin t>0$: dla $t\le a$ wynika to z poprzedniego dowodu, a dla $a<t<2a$ z $\sin(\pi-t)=\sin t$. Dla $0\le u<v\le\pi$ wzory dodawania dają
$\cos u-\cos v=2\sin((u+v)/2)\sin((v-u)/2)>0$.
Dla $-a\le u<v\le a$ analogicznie
$\sin v-\sin u=2\cos((u+v)/2)\sin((v-u)/2)>0$, ponieważ $|(u+v)/2|<a$. Zakresy wartości wynikają z wartości końcowych i twierdzenia o wartości pośredniej.

Z $\cos(\pi+t)=-\cos t$ i $\sin(\pi+t)=-\sin t$ wynika, że sinus jest ujemny na $(\pi,2\pi)$. Jego jedyne zera na $[0,2\pi)$ to $0,\pi$, a $\cos\pi=-1$. Jeżeli $T>0$ jest okresem cosinusa, to $\cos T=1$, zatem $\sin T=0$ i po odjęciu całkowitej wielokrotności $2\pi$ pozostaje tylko reszta $0$. Jeżeli $T>0$ jest okresem sinusa, to $\sin T=0$ oraz $\sin(a+T)=\sin a=1$, czyli z wzoru dodawania $\cos T=1$. Ten sam argument daje tezę.
\end{proof}
\begin{center}
\begin{tikzpicture}[x=1.65cm,y=1.5cm,>=Stealth]
 \draw[->] (-3.45,0)--(3.55,0) node[right] {$t$};
 \draw[->] (0,-1.3)--(0,1.4) node[above] {$y$};
 \draw[blue!65!black,thick,domain=-3.14:3.14,samples=100] plot (\x,{sin(\x r)});
 \draw[orange!80!black,thick,dashed,domain=-3.14:3.14,samples=100] plot (\x,{cos(\x r)});
 \foreach \x/\lab in {-3.14159/{-\pi},-1.5708/{-\pi/2},1.5708/{\pi/2},3.14159/{\pi}}{
  \draw (\x,.05)--(\x,-.05) node[below] {$\lab$};}
 \node[blue!65!black] at (2.2,1.17) {$\sin t$};
 \node[orange!80!black] at (-2.4,-1.16) {$\cos t$};
\end{tikzpicture}

\small Rysunek: sinus i cosinus; etykiety argumentów podano w radianach.
\end{center}

\section{Tangens i funkcje cyklometryczne}
\begin{definicja}
\label{ma1:def:inverse-trig}
Definiujemy $\tan t=\sin t/\cos t$ tam, gdzie $\cos t\ne0$, oraz $\cot t=\cos t/\sin t$ tam, gdzie $\sin t\ne0$. Funkcje $\arcsin$ i $\arccos$ są odwrotnościami ograniczeń sinusa do $[-\pi/2,\pi/2]$ i cosinusa do $[0,\pi]$.
\end{definicja}
Z twierdzenia \ref{ma1:thm:trig-monotonicity} oraz wzorów przesunięcia o $\pi$ wynika, że dziedziny tangensa i cotangensa to odpowiednio $\R\setminus\{\pi/2+k\pi:k\in\Z\}$ oraz $\R\setminus\{k\pi:k\in\Z\}$. Są ciągłe na swoich dziedzinach i mają okres $\pi$. Funkcje $\arcsin,\arccos:[-1,1]\to\R$ są ciągłe dzięki twierdzeniu \ref{ma1:thm:strict-monotone-inverse}; pierwsza jest ściśle rosnąca, druga ściśle malejąca.
\begin{twierdzenie}
\label{ma1:thm:tan-inverse}
Tangens na $(-\pi/2,\pi/2)$ jest ciągłą ściśle rosnącą bijekcją na $\R$, a cotangens na $(0,\pi)$ jest ciągłą ściśle malejącą bijekcją na $\R$. Ich odwrotności oznaczamy odpowiednio przez $\arctan$ i $\operatorname{arccot}$; są ciągłe.
\end{twierdzenie}
\begin{proof}
Dla $-a<u<v<a$ mamy
$\tan v-\tan u=\sin(v-u)/(\cos v\cos u)>0$, gdyż $0<v-u<\pi$ i oba cosinusy są dodatnie. Przy $t\to a^-$ sinus dąży do $1$, a dodatni cosinus do $0$, zatem iloraz dąży do $+\infty$: dla $M>0$ wystarczy $\sin t>1/2$ i $0<\cos t<1/(2M)$. Nieparzystość daje granicę $-\infty$ przy $-a$ z prawej. Twierdzenie o wartości pośredniej i ścisła monotoniczność dają bijektywność. Ponieważ $\cot t=\tan(a-t)$ dla $0<t<\pi$, otrzymujemy analogiczną tezę z przeciwnym kierunkiem monotoniczności. Ciągłość odwrotności wynika z twierdzenia \ref{ma1:thm:strict-monotone-inverse}.
\end{proof}
W szczególności $\arccos x=\pi/2-\arcsin x$ dla $|x|\le1$ oraz $\operatorname{arccot}x=\pi/2-\arctan x$ dla $x\in\R$. Pierwsza równość wynika z $\cos(a-t)=\sin t$ oraz jednoznaczności odwrotności, druga z analogicznej tożsamości użytej w dowodzie.

\section{Funkcje hiperboliczne}
\begin{definicja}
\label{ma1:def:hyperbolic}
Dla $x\in\R$ przyjmujemy
$\sinh x=(\exp x-\exp(-x))/2$,
$\cosh x=(\exp x+\exp(-x))/2$ oraz $\tanh x=\sinh x/\cosh x$. Funkcję $\coth x=\cosh x/\sinh x$ określamy dla $x\ne0$.
\end{definicja}
Z prawa wykładniczego wynika $\cosh^2x-\sinh^2x=1$, ponieważ różnica kwadratów jest równa $\exp x\exp(-x)$. Funkcja $\sinh$ jest nieparzysta, $\cosh$ parzysta, a $\cosh x\ge1$ (jest dodatnia i jej kwadrat jest nie mniejszy niż $1$). Rozwijając iloczyny i grupując składniki według sum oraz różnic argumentów, otrzymujemy
\begin{align*}
 \cosh(x+y)&=\cosh x\cosh y+\sinh x\sinh y,\\
 \sinh(x+y)&=\sinh x\cosh y+\cosh x\sinh y.
\end{align*}
Istotnie, po podstawieniu definicji po prawej stronie pierwszego wzoru pozostają składniki $\exp(x+y)/2$ i $\exp(-x-y)/2$, a w drugim ich różnica.

\begin{twierdzenie}
\label{ma1:thm:hyperbolic-inverse}
Funkcja $\sinh:\R\to\R$ jest ściśle rosnącą bijekcją, $\cosh:[0,+\infty)\to[1,+\infty)$ jest ściśle rosnącą bijekcją, a $\tanh:\R\to(-1,1)$ jest ściśle rosnącą bijekcją. Ich ciągłe odwrotności mają wzory
\begin{align*}
 \operatorname{arsinh}y&=\log(y+\sqrt{y^2+1}), &&y\in\R,\\
 \operatorname{arcosh}y&=\log(y+\sqrt{y^2-1}), &&y\ge1,\\
 \operatorname{artanh}y&=\tfrac12\log\frac{1+y}{1-y}, &&|y|<1.
\end{align*}
\end{twierdzenie}
\begin{proof}
Jeśli $x<z$, to $\exp x<\exp z$ i $\exp(-x)>\exp(-z)$, więc $\sinh x<\sinh z$. Granice w nieskończonościach wynoszą odpowiednio $+\infty$ i $-\infty$ z definicji oraz nieparzystości. Na $[0,+\infty)$ funkcja $\sinh$ jest nieujemna i ściśle rosnąca, dlatego $\cosh=\sqrt{1+\sinh^2}$ jest ściśle rosnąca, zaczyna się od $1$ i jest nieograniczona. Ciągłość i twierdzenie o wartości pośredniej dają oba obrazy.

Mamy $\tanh x=(\exp(2x)-1)/(\exp(2x)+1)$. Funkcja $t\mapsto(t-1)/(t+1)=1-2/(t+1)$ jest ściśle rosnąca dla $t>0$ i przyjmuje dokładnie wartości z $(-1,1)$, ponieważ dla $|y|<1$ równanie $y=(t-1)/(t+1)$ ma dodatnie rozwiązanie $t=(1+y)/(1-y)$. To dowodzi bijektywności i wzoru dla $\operatorname{artanh}$.

Podstawmy $t=\exp x>0$ w równaniu $\sinh x=y$. Otrzymujemy $t^2-2yt-1=0$, którego jedynym dodatnim pierwiastkiem jest $y+\sqrt{y^2+1}$. Dla $\cosh x=y$, $x\ge0$, mamy $t\ge1$ i $t^2-2yt+1=0$. Dla $y>1$ pierwiastek $y+\sqrt{y^2-1}>1$, a drugi jest jego odwrotnością, mniejszą od $1$; dla $y=1$ oba wynoszą $1$. Stosujemy $\log$ i otrzymujemy podane wzory. Ciągłość wszystkich odwrotności wynika z twierdzenia \ref{ma1:thm:strict-monotone-inverse}.
\end{proof}

\section{Podstawowe granice}
\begin{twierdzenie}
\label{ma1:thm:elementary-limits}
Zachodzą granice
\[
 \lim_{t\to0}\frac{\exp t-1}{t}=1,\qquad
 \lim_{t\to0}\frac{\sin t}{t}=1,\qquad
 \lim_{t\to0}\frac{1-\cos t}{t^2}=\frac12.
\]
Ponadto $\lim_{u\to0}\log(1+u)/u=1$ i dla $\alpha\in\R$,
$\lim_{u\to0}((1+u)^\alpha-1)/u=\alpha$.
\end{twierdzenie}
\begin{proof}
Dla $0<|t|\le1$ mamy
$|(\exp t-1)/t-1|\le |t|\sum_{n=2}^\infty1/n!$.
Podobnie $|\sin t/t-1|\le t^2\sum_{n=1}^\infty1/(2n+1)!$.
Obie prawe strony dążą do zera. Trzecia granica wynika z $1-\cos t=2\sin^2(t/2)$, a więc iloraz jest równy $\tfrac12(\sin(t/2)/(t/2))^2$.

Dla $u>-1$, $u\ne0$, podstawmy $v=\log(1+u)$. Wówczas $v\ne0$, $v\to0$ przy $u\to0$ z ciągłości logarytmu, a $u=\exp v-1$. Stąd $\log(1+u)/u=v/(\exp v-1)\to1$ z pierwszej granicy. Dla $\alpha=0$ ostatni iloraz jest zerowy; dla $\alpha\ne0$ zapisujemy
\[
 \frac{(1+u)^\alpha-1}{u}
 =\frac{\exp(\alpha v)-1}{\alpha v}\cdot\alpha\frac vu.
\]
Obie granice czynników dają $\alpha$.
\end{proof}

\section{Postać trygonometryczna liczb zespolonych}
\begin{twierdzenie}
\label{ma1:thm:polar-complex}
Każde $z\in\C\setminus\{0\}$ można zapisać jako $z=r(\cos t+i\sin t)$, gdzie $r=|z|>0$, a $t$ jest jednoznaczny z dokładnością do wielokrotności $2\pi$. Dla $n\in\Z$ zachodzi wzór de Moivre'a
$z^n=r^n(\cos(nt)+i\sin(nt))$.
\end{twierdzenie}
\begin{proof}
Niech $u=\operatorname{Re}z/r$, $v=\operatorname{Im}z/r$. Mamy $u^2+v^2=1$. Jeśli $v\ge0$, wybieramy $t=\arccos u\in[0,\pi]$; wtedy $\sin t\ge0$ i $\sin^2t=v^2$, więc $\sin t=v$. Jeśli $v<0$, wybieramy $t=-\arccos u$, co daje ujemny sinus i tę samą wartość cosinusa. Jeśli dwa argumenty $t,s$ dają ten sam punkt okręgu, to $\exp(i(t-s))=1$; z opisu zer sinusa i wartości cosinusa w poprzednim twierdzeniu wynika $t-s\in2\pi\Z$. Odwrotna implikacja wynika z okresowości. Wzór de Moivre'a jest konsekwencją prawa wykładniczego stosowanego indukcyjnie dla $n\ge0$ oraz $\exp(-it)=1/\exp(it)$ dla $n<0$.
\end{proof}
\begin{wniosek}
\label{ma1:cor:complex-nth-roots}
Dla $w=r(\cos t+i\sin t)\ne0$ i całkowitego $n\ge1$ równanie $z^n=w$ ma dokładnie $n$ różnych rozwiązań:
\[
 z_k=\sqrt[n]r\left(\cos\frac{t+2k\pi}{n}
                  +i\sin\frac{t+2k\pi}{n}\right),
 \qquad k=0,\ldots,n-1.
\]
Dla $w=0$ jedynym rozwiązaniem jest $z=0$.
\end{wniosek}
\begin{proof}
Wzór de Moivre'a dowodzi, że podane liczby są rozwiązaniami. Dla dowolnego rozwiązania niezerowego postać trygonometryczna i multiplikatywność modułu dają $|z|=\sqrt[n]r$, a argument $s$ spełnia $ns-t\in2\pi\Z$. Redukując odpowiednią liczbę całkowitą modulo $n$, otrzymujemy jedno z wypisanych rozwiązań. Dwa różne indeksy z zakresu $0,\ldots,n-1$ dają różnicę argumentów, która nie jest wielokrotnością $2\pi$, więc rozwiązania są różne. Jeśli $z^n=0$, niezerowość iloczynu niezerowych liczb w ciele $\C$ wymusza $z=0$.
\end{proof}
